\chapter{Kissing Disease (Glandular Fever)}

\section*{Definition and Overview}
"Kissing disease" is the popular name for glandular fever, also known as infectious mononucleosis, an illness caused by the Epstein--Barr virus (EBV). The name reflects the way the virus is commonly spread through saliva, including by kissing. Glandular fever typically affects teenagers and young adults and is characterised by a triad of fever, severe sore throat, and swollen lymph glands. While the illness can be unpleasant and the fatigue prolonged, most people recover fully with rest and supportive care. This chapter explains the causes, symptoms, and management of glandular fever and its potential complications.

\section*{Epidemiology}
EBV infection is extremely common worldwide. Most people are infected by early adulthood, and in childhood the infection is often mild or produces no symptoms at all. When infection is delayed until adolescence or young adulthood, as occurs in many higher-income countries, a symptomatic illness with the features of glandular fever develops in a substantial proportion of cases. Glandular fever occurs in people of all ages but is most common between 15 and 24 years. Outbreaks are frequent in settings where young people live closely together, such as universities and boarding schools. Once infected, people carry the virus for life in a dormant state, but the illness itself usually occurs only once.

\section*{Aetiology and Pathophysiology}
\begin{itemize}
    \item \textbf{Cause:} the Epstein--Barr virus, a member of the herpesvirus family, is the causative agent
    \item \textbf{Transmission:} spread through saliva, via kissing, sharing drinks or utensils, or coughing and sneezing
    \item \textbf{Incubation:} symptoms typically begin four to six weeks after exposure
    \item \textbf{Pathophysiology:} the virus infects B lymphocytes, which mount a vigorous immune response dominated by atypical T lymphocytes, producing the characteristic swollen lymph nodes, enlarged spleen, and sore throat
    \item \textbf{Latency:} after the acute illness the virus remains dormant in B cells for life and can reactivate silently, allowing occasional shedding in saliva
\end{itemize}

\section*{Clinical Features}
Symptoms typically develop gradually over a few days and last two to four weeks, though fatigue may persist for months.
\begin{itemize}
    \item Severe sore throat, sometimes with difficulty swallowing
    \item Fever, often with sweating and chills
    \item Swollen, tender lymph nodes in the neck, armpits, and groin
    \item Extreme tiredness and fatigue, sometimes lasting weeks or months
    \item Enlargement of the spleen and sometimes the liver
    \item Headache, muscle aches, and loss of appetite
    \item Swollen eyelids or a fine pink rash, especially if amoxicillin or ampicillin antibiotics are taken
\end{itemize}

\section*{Differential Diagnosis}
\begin{itemize}
    \item \textbf{Streptococcal tonsillitis:} bacterial sore throat with fever and swollen glands, distinguished by throat swab
    \item \textbf{Other viral infections:} cytomegalovirus, HIV seroconversion illness, and adenovirus can mimic glandular fever
    \item \textbf{Toxoplasmosis:} causes fever, lymphadenopathy, and fatigue
    \item \textbf{Acute leukaemia or lymphoma:} rare causes of persistent fever, weight loss, and lymphadenopathy
    \item \textbf{Medication reactions} and autoimmune conditions causing lymph node swelling
\end{itemize}

\section*{When to Seek Medical Care}
Most cases can be managed at home, but medical review is recommended to confirm the diagnosis, particularly when symptoms are severe or prolonged. See a doctor if the sore throat prevents drinking, symptoms worsen after initial improvement, or fever persists beyond a week.

\textbf{Call 000 (or local emergency services) for:}
\begin{itemize}
    \item Severe abdominal pain, especially in the upper left side, which may indicate a ruptured spleen
    \item Difficulty breathing or swallowing, or drooling
    \item Confusion, severe headache, stiff neck, or seizures
    \item Very pale or yellow skin, or unexplained bruising or bleeding
\end{itemize}

\section*{Diagnosis and Investigations}
\begin{itemize}
    \item \textbf{History and examination:} the classic triad of fever, sore throat, and lymphadenopathy, particularly in a young person, strongly suggests the diagnosis
    \item \textbf{Blood tests:} the Monospot test or EBV antibody tests confirm recent infection
    \item \textbf{Full blood count:} shows characteristic atypical lymphocytes; also used to exclude leukaemia if the presentation is unusual
    \item \textbf{Liver function tests:} mildly elevated liver enzymes are common and usually resolve spontaneously
    \item \textbf{Throat swab:} to exclude concurrent streptococcal infection if suspected
\end{itemize}

\section*{Management}
\begin{itemize}
    \item \textbf{Rest:} adequate rest during the acute illness, then a gradual return to normal activity
    \item \textbf{Pain and fever relief:} paracetamol or ibuprofen for sore throat, fever, and aches
    \item \textbf{Fluids:} generous fluids, cool drinks, and soft foods to ease swallowing
    \item \textbf{Gargles and lozenges:} salt-water gargles or throat lozenges may soothe the throat
    \item \textbf{Avoid contact sport:} for at least four weeks, or longer if the spleen remains enlarged, to reduce the risk of splenic rupture
    \item \textbf{Antibiotics:} not effective against the virus; avoid amoxicillin and ampicillin because they frequently cause a rash in glandular fever
    \item \textbf{Corticosteroids:} occasionally used for severe throat swelling that obstructs the airway, under specialist guidance
\end{itemize}

\section*{Complications}
\begin{itemize}
    \item \textbf{Splenic rupture:} rare but serious, causing sudden severe abdominal pain and collapse; a surgical emergency
    \item Airway obstruction from severe tonsillar swelling, in rare cases
    \item Hepatitis with jaundice, usually mild and self-limiting
    \item Anaemia, low platelet count, or rarely haemophagocytic lymphohistiocytosis
    \item Chronic fatigue that persists for months in some people
    \item Neurological complications are rare, including meningitis, encephalitis, or Guillain--Barré syndrome
\end{itemize}

\section*{Prognosis and Outcomes}
Glandular fever is almost always a self-limiting illness. The acute symptoms settle over two to four weeks, and most people feel back to normal within two months. Fatigue resolves more slowly in some people, but persistent symptoms beyond six months are uncommon and warrant medical review. Full recovery is the rule, and once recovered, people rarely experience the illness again, although the virus remains dormant in the body for life.

\section*{Prevention and Self-Care}
\begin{itemize}
    \item Avoid sharing drinks, utensils, or toothbrushes with someone who has the illness
    \item Avoid kissing when symptoms are present, since the virus is shed in saliva
    \item Practise good hand hygiene, especially after contact with respiratory secretions
    \item Rest and stay well hydrated during the illness
    \item Take time off school or work while feverish, and allow the body to recover gradually
    \item Avoid vigorous exercise and contact sport until cleared by a doctor
\end{itemize}

\section*{Sources and Further Reading}
The following reputable sources provide further information for healthcare students and patients seeking reliable, current advice about glandular fever.
\begin{itemize}
    \item Healthdirect. \textit{Glandular fever.} https://www.healthdirect.gov.au/glandular-fever
    \item NHS. \textit{Glandular fever (infectious mononucleosis).} https://www.nhs.uk/conditions/glandular-fever/
    \item Better Health Channel (Victoria). \textit{Glandular fever.} https://www.betterhealth.vic.gov.au/
    \item Mayo Clinic. \textit{Infectious mononucleosis.} https://www.mayoclinic.org/diseases-conditions/mononucleosis/
    \item MSD Manual. \textit{Infectious mononucleosis.} https://www.msdmanuals.com/
    \item Centers for Disease Control and Prevention. \textit{Epstein--Barr virus and infectious mononucleosis.} https://www.cdc.gov/epstein-barr/
\end{itemize}
